\hmsubsection{Tools and Workflow}{OAH}{UMA}\label{sec:tools-workflow}

This section describes the tools and platforms used to support the team's
development process, communication, and documentation throughout the project.

\textbf{Project management.} Task planning and tracking were managed using a
digital Trello board, organised into four lanes: Backlog, In Progress, Testing,
and Done. Tasks were broken down into actionable items at the start of each sprint,
and ownership was assigned per member to maintain clear accountability.

\textbf{Communication.} Internal communication was handled through a dedicated
Microsoft Teams channel and Slack, which supported both asynchronous messaging and
video meetings for members unable to attend in person. Shared documents, steering
documents, and time registration were stored on a common OneDrive accessible to
all members.

\textbf{Version control.} Source code was managed using Git with a shared GitHub
repository. The team committed daily, ensuring that all members had access to the
latest version of the codebase at all times.

\textbf{Development environment.} Host-side Python code was developed in Visual
Studio Code. Firmware for the OpenRB-150 motor controller was written and flashed
using the Arduino IDE~2.x. The full technology stack is described in
Section~\ref{sec:technology-stack} and summarised in Table~\ref{tab:tech-stack}.

\textbf{Documentation.} The project report was written collaboratively in
\LaTeX{} using Overleaf, which allowed all members to contribute simultaneously
without version conflicts.
